% !TEX program = xelatex
% ============================================================
%  متتبع المعالم — Milestone Tracker Template
%  نموذج لتتبع معالم المشروع مع مخطط زمني تفاعلي
%  Compile with: xelatex main.tex (run twice)
% ============================================================
%  Features:
%    - TikZ timeline diagram showing milestones on a horizontal axis
%    - Milestone table: ID, Name, Planned date, Actual date,
%      Status, Dependencies, Notes
%    - Status summary section
% ============================================================

\documentclass[11pt,a4paper]{article}

\usepackage[a4paper,margin=2.5cm,top=2.5cm,bottom=2.5cm,headheight=16pt]{geometry}
\usepackage{graphicx}
\usepackage{array}
\usepackage{booktabs}
\usepackage{tabularx}
\usepackage{multirow}
\usepackage{enumitem}
\usepackage{float}
\usepackage{titlesec}
\usepackage{fancyhdr}
\usepackage{setspace}
\usepackage{xcolor}
\usepackage{amsmath}
\usepackage{tikz}
\usetikzlibrary{positioning,calc,shapes.geometric,decorations.pathreplacing,arrows.meta}

\usepackage{bidi}
\usepackage[no-math]{fontspec}
\usepackage[quiet]{polyglossia}

\setmainfont[Scale=1]{Amiri-Regular.ttf}
\newfontfamily\arabicfont[Script=Arabic,Scale=0.95]{Amiri-Regular.ttf}
\newfontfamily\englishfont[Scale=0.95]{TeX Gyre Termes}

\setdefaultlanguage[calendar=gregorian,hijricorrection=1]{arabic}
\setotherlanguage[variant=british]{english}

\usepackage[breaklinks,hidelinks]{hyperref}

\addto\captionsarabic{%
  \renewcommand{\contentsname}{الفهرس}%
  \renewcommand{\figurename}{الشكل}%
  \renewcommand{\tablename}{الجدول}%
}

\titleformat{\section}{\Large\bfseries}{\thesection}{0.7em}{}
\titlespacing*{\section}{0pt}{1.4ex plus 0.4ex}{0.9ex plus 0.3ex}

\pagestyle{fancy}
\fancyhf{}
\fancyhead[R]{\small\itshape متتبع المعالم}
\fancyhead[L]{\small اسم المشروع}
\fancyfoot[C]{\small\thepage}
\renewcommand{\headrulewidth}{0.3pt}

\setlist[itemize]{topsep=0.3em, itemsep=0.15em, parsep=0pt}
\setlist[enumerate]{topsep=0.3em, itemsep=0.15em, parsep=0pt}

\newcolumntype{L}[1]{>{\raggedright\arraybackslash}p{#1}}
\newcolumntype{C}[1]{>{\centering\arraybackslash}p{#1}}

\definecolor{msDone}{HTML}{27AE60}
\definecolor{msProgress}{HTML}{F39C12}
\definecolor{msPending}{HTML}{95A5A6}
\definecolor{msAxis}{HTML}{1A3C6B}

\setstretch{1.3}

\begin{document}

% ============================================================
% TITLE BLOCK
% ============================================================
\begin{center}
{\LARGE\bfseries متتبع المعالم}\\[0.3em]
{\large Milestone Tracker}\\[1em]
\end{center}

% ============================================================
% PROJECT INFO
% ============================================================
% TODO: املأ بيانات المشروع.
\begin{table}[H]
\centering
\renewcommand{\arraystretch}{1.4}
\begin{tabularx}{\textwidth}{L{3.5cm} X L{3.5cm} X}
\toprule
\textbf{اسم المشروع:} & --- & \textbf{مدير المشروع:} & --- \\
\textbf{اسم الشركة:} & --- & \textbf{تاريخ التحديث:} & --- \\
\textbf{تاريخ البدء:} & --- & \textbf{تاريخ الانتهاء:} & --- \\
\bottomrule
\end{tabularx}
\end{table}

\vspace{1em}

% ============================================================
% 1. TIMELINE DIAGRAM
% ============================================================
\section{المخطط الزمني للمعالم}
\label{sec:timeline}

% TODO: عدّل مواقع المعالم على المخطط حسب تواريخ مشروعك.
% The timeline uses LTR coordinates; Arabic labels are placed via \RL{}.
\begin{figure}[H]
\centering
\begin{tikzpicture}[x=1cm,y=1cm]

% ---- Axis ----
\draw[msAxis,very thick,-{Stealth[length=3mm]}] (0,0) -- (13,0);

% ---- Month markers ----
\foreach \x/\label in {0/الشهر 1, 2/الشهر 3, 4/الشهر 5, 6/الشهر 7, 8/الشهر 9, 10/الشهر 11, 12/الشهر 12} {
  \draw[msAxis,thick] (\x,-0.15) -- (\x,0.15);
  \node[below,font=\scriptsize,msAxis] at (\x,-0.2) {\label};
}

% ---- Milestones ----
% M1: done (month 2)
\node[diamond,fill=msDone,draw=msDone,inner sep=4pt] (m1) at (1,1.2) {};
\node[above=2pt of m1,font=\scriptsize\bfseries] {\RL{المعلم 1}};
\node[above=14pt of m1,font=\scriptsize] {\RL{بدء المشروع}};

% M2: done (month 4)
\node[diamond,fill=msDone,draw=msDone,inner sep=4pt] (m2) at (3,1.2) {};
\node[above=2pt of m2,font=\scriptsize\bfseries] {\RL{المعلم 2}};
\node[above=14pt of m2,font=\scriptsize] {\RL{انتهاء التصميم}};

% M3: in progress (month 6)
\node[diamond,fill=msProgress,draw=msProgress,inner sep=4pt] (m3) at (5,1.2) {};
\node[above=2pt of m3,font=\scriptsize\bfseries] {\RL{المعلم 3}};
\node[above=14pt of m3,font=\scriptsize] {\RL{انتهاء التطوير}};

% M4: pending (month 8)
\node[diamond,fill=msPending,draw=msPending,inner sep=4pt] (m4) at (7,1.2) {};
\node[above=2pt of m4,font=\scriptsize\bfseries] {\RL{المعلم 4}};
\node[above=14pt of m4,font=\scriptsize] {\RL{انتهاء الاختبار}};

% M5: pending (month 10)
\node[diamond,fill=msPending,draw=msPending,inner sep=4pt] (m5) at (9,1.2) {};
\node[above=2pt of m5,font=\scriptsize\bfseries] {\RL{المعلم 5}};
\node[above=14pt of m5,font=\scriptsize] {\RL{الإطلاق التجريبي}};

% M6: pending (month 12)
\node[diamond,fill=msPending,draw=msPending,inner sep=4pt] (m6) at (11,1.2) {};
\node[above=2pt of m6,font=\scriptsize\bfseries] {\RL{المعلم 6}};
\node[above=14pt of m6,font=\scriptsize] {\RL{الإطلاق النهائي}};

% ---- Vertical dashed lines from milestones to axis ----
\foreach \m in {m1,m2,m3,m4,m5,m6} {
  \draw[dashed,gray,thin] (\m) -- ($(\m |- 0,0)$);
}

% ---- Legend ----
\node[diamond,fill=msDone,draw=msDone,inner sep=3pt] at (1.5,-1.5) {};
\node[right=2pt,font=\scriptsize] at (1.7,-1.5) {\RL{مكتمل}};

\node[diamond,fill=msProgress,draw=msProgress,inner sep=3pt] at (4.5,-1.5) {};
\node[right=2pt,font=\scriptsize] at (4.7,-1.5) {\RL{قيد التنفيذ}};

\node[diamond,fill=msPending,draw=msPending,inner sep=3pt] at (8,-1.5) {};
\node[right=2pt,font=\scriptsize] at (8.2,-1.5) {\RL{لم يبدأ}};

\end{tikzpicture}
\caption{المخطط الزمني لمعالم المشروع}
\end{figure}

% ============================================================
% 2. MILESTONE TABLE
% ============================================================
\section{جدول المعالم}
\label{sec:milestones}

% TODO: عدّل صفوف المعالم حسب مشروعك.
\begin{table}[H]
\centering
\caption{تفاصيل معالم المشروع}
\renewcommand{\arraystretch}{1.4}
\small
\begin{tabularx}{\textwidth}{C{1.2cm} L{3cm} L{2cm} L{2cm} L{2cm} L{2.5cm} X}
\toprule
\textbf{المعرف} & \textbf{اسم المعلم} & \textbf{التاريخ المخطط} & \textbf{التاريخ الفعلي} & \textbf{الحالة} & \textbf{التبعيات} & \textbf{ملاحظات} \\
\midrule
\LR{M1} & بدء المشروع & --- & --- & مكتمل & --- & تم البدء في الموعد المحدد \\
\LR{M2} & انتهاء التصميم & --- & --- & مكتمل & \LR{M1} & اكتمل التصميم بنجاح \\
\LR{M3} & انتهاء التطوير & --- & --- & قيد التنفيذ & \LR{M2} & تأخر بسيط بسبب متطلبات إضافية \\
\LR{M4} & انتهاء الاختبار & --- & --- & لم يبدأ & \LR{M3} & يبدأ بعد اكتمال التطوير \\
\LR{M5} & الإطلاق التجريبي & --- & --- & لم يبدأ & \LR{M4} & يتطلب موافقة لجنة الجودة \\
\LR{M6} & الإطلاق النهائي & --- & --- & لم يبدأ & \LR{M5} & الإطلاق الرسمي للمستخدمين \\
\bottomrule
\end{tabularx}
\end{table}

% ============================================================
% 3. STATUS SUMMARY
% ============================================================
\section{ملخص الحالة}
\label{sec:summary}

% TODO: حدّث ملخص الحالة حسب التقدم الفعلي.
\begin{table}[H]
\centering
\renewcommand{\arraystretch}{1.4}
\begin{tabularx}{\textwidth}{L{5cm} C{3cm} X}
\toprule
\textbf{المؤشر} & \textbf{القيمة} & \textbf{ملاحظات} \\
\midrule
إجمالي المعالم & 6 & --- \\
المعالم المكتملة & 2 & \LR{M1, M2} \\
المعالم قيد التنفيذ & 1 & \LR{M3} \\
المعالم المتأخرة & 0 & --- \\
المعالم المتبقية & 3 & \LR{M4, M5, M6} \\
نسبة الإنجاز الكلية & 33\% & تقديري بناءً على المعالم المكتملة \\
\bottomrule
\end{tabularx}
\end{table}

\end{document}
